\documentclass[12pt,a4paper]{article}
\usepackage[utf8]{inputenc}
\usepackage[margin=1in]{geometry}
\usepackage{graphicx}
\usepackage{hyperref}
\usepackage{enumitem}
\usepackage{tabularx}
\usepackage{booktabs}
\usepackage{float}
\usepackage{amsmath}
\usepackage{xcolor}
\usepackage{titlesec}
\usepackage{fancyhdr}
\usepackage{minted}
\usepackage{longtable}
\usepackage{multirow}

% Page formatting
\pagestyle{fancy}
\fancyhf{}
\fancyhead[L]{School Management Application - ITPM Project Report}
\fancyhead[R]{\thepage}
\fancyfoot[C]{Confidential}

% Title formatting
\titleformat{\section}{\Large\bfseries}{\thesection}{1em}{}
\titleformat{\subsection}{\large\bfseries}{\thesubsection}{1em}{}
\titleformat{\subsubsection}{\normalsize\bfseries}{\thesubsubsection}{1em}{}

% Minted settings
\setminted{
    fontsize=\small,
    breaklines=true,
    frame=single,
    framesep=2mm,
    baselinestretch=1.1,
    bgcolor=gray!5
}

\title{\textbf{School Management Application (SMA)}\\
       \large IT Project Management - Final Project Report\\[1em]
       \normalsize Agile Scrum Methodology Implementation}
\author{Project Team: 5 Members\\
        Course: IT Project Management\\
        Submission Date: \today}
\date{}

\begin{document}

\maketitle
\thispagestyle{empty}
\newpage

\tableofcontents
\newpage

%==============================================================================
\section{Executive Summary}
%==============================================================================

This report documents the comprehensive development of a School Management Application (SMA) designed to streamline administrative operations, student management, and staff coordination in educational institutions. The project was developed using modern web technologies and follows Object-Oriented Programming (OOP) principles with a Service-Oriented Architecture.

The application is built using \textbf{Next.js 16} with React 19, utilizing \textbf{Supabase} as a Backend-as-a-Service (BaaS) platform providing PostgreSQL database and authentication infrastructure. The development followed \textbf{Agile Scrum methodology} with a team of 5 members across 4 sprints spanning 12 weeks.

\textbf{Key Achievements:}
\begin{itemize}[noitemsep]
    \item Implemented secure JWT-based authentication with role-based access control (RBAC)
    \item Developed comprehensive Student Portal with dashboard, fees, attendance, marks, and exam schedules
    \item Created Staff Portal with salary management, class management, and attendance tracking
    \item Built Admin Panel for system-wide management and feedback handling
    \item Established a robust service layer architecture following OOP principles
    \item Deployed production-ready application with responsive design and dark mode support
\end{itemize}

\textbf{Note:} The originally planned ``Assign Room'' feature was removed from scope during Sprint 3 due to complexity constraints and timeline considerations.

%==============================================================================
\section{Project Overview}
%==============================================================================

\subsection{Project Background}

Educational institutions face significant challenges in managing student records, staff information, attendance tracking, and academic performance data. Traditional paper-based systems are inefficient, error-prone, and difficult to maintain. The School Management Application (SMA) addresses these challenges by providing a centralized, web-based platform for all stakeholders.

\subsection{Project Objectives}

The primary objectives of this project are:

\begin{enumerate}[noitemsep]
    \item Develop a secure, scalable web-based school management system
    \item Implement role-based access control for Students, Staff, and Administrators
    \item Provide real-time dashboards with key performance indicators
    \item Enable efficient attendance tracking and grade management
    \item Create intuitive fee and salary management modules
    \item Support feedback collection and management system
    \item Ensure responsive design for multiple device types
    \item Follow software engineering best practices with OOP architecture
\end{enumerate}

\subsection{Project Scope}

\subsubsection{In Scope}
\begin{itemize}[noitemsep]
    \item \textbf{Authentication System:} Secure login, logout, session management, password change
    \item \textbf{Student Portal:} Dashboard, fee management, exam schedule, timetable, attendance viewing, marks viewing, profile management, feedback submission
    \item \textbf{Staff Portal:} Dashboard, salary viewing, class management, attendance marking, timetable management, profile management
    \item \textbf{Admin Panel:} Dashboard statistics, student management, staff management, class management, feedback management, user administration
    \item \textbf{Database:} PostgreSQL via Supabase with 13 interconnected tables
    \item \textbf{API Layer:} RESTful API endpoints using Next.js App Router
    \item \textbf{Security:} JWT authentication, password hashing with bcrypt, RBAC
\end{itemize}

\subsubsection{Out of Scope}
\begin{itemize}[noitemsep]
    \item Mobile application (iOS/Android)
    \item Parent portal
    \item Payment gateway integration
    \item Email/SMS notification system
    \item \textbf{Room assignment feature (dropped during development)}
\end{itemize}

\subsection{Stakeholders}

\begin{table}[H]
\centering
\begin{tabularx}{\textwidth}{|l|X|c|}
\hline
\textbf{Stakeholder} & \textbf{Role \& Interest} & \textbf{Priority} \\
\hline
Students & Primary end users; access academic info, fees, schedules & High \\
\hline
Teaching Staff & Manage classes, mark attendance, view salary & High \\
\hline
School Admin & System oversight, user management, reports & High \\
\hline
Development Team & Design, development, testing, deployment & High \\
\hline
Project Manager & Coordination, timeline management, risk mitigation & High \\
\hline
IT Department & Infrastructure, security, maintenance & Medium \\
\hline
\end{tabularx}
\caption{Project Stakeholders Matrix}
\end{table}

%==============================================================================
\section{Agile Methodology \& Scrum Framework}
%==============================================================================

\subsection{Why Agile?}

The project adopted Agile methodology for several strategic reasons:

\begin{enumerate}[noitemsep]
    \item \textbf{Iterative Development:} Allows for continuous improvement and early feedback
    \item \textbf{Flexibility:} Accommodates changing requirements (e.g., dropping room assignment feature)
    \item \textbf{Risk Mitigation:} Early identification of issues through frequent deliverables
    \item \textbf{Stakeholder Engagement:} Regular demos and reviews keep stakeholders informed
    \item \textbf{Team Collaboration:} Daily standups promote communication and problem-solving
\end{enumerate}

\subsection{Scrum Framework Implementation}

\subsubsection{Sprint Structure}
\begin{itemize}[noitemsep]
    \item \textbf{Sprint Duration:} 3 weeks per sprint
    \item \textbf{Total Sprints:} 4 sprints (12 weeks total)
    \item \textbf{Sprint Planning:} 3 hours at sprint start
    \item \textbf{Daily Standup:} 15 minutes each morning
    \item \textbf{Sprint Review:} 2 hours at sprint end
    \item \textbf{Sprint Retrospective:} 1.5 hours after review
\end{itemize}

\subsubsection{Scrum Artifacts}
\begin{itemize}[noitemsep]
    \item \textbf{Product Backlog:} Prioritized list of all features and requirements
    \item \textbf{Sprint Backlog:} Selected items for current sprint
    \item \textbf{Increment:} Potentially shippable product at sprint end
    \item \textbf{Definition of Done (DoD):} Code reviewed, tested, documented, deployed to staging
\end{itemize}

\subsubsection{Scrum Events}
\begin{itemize}[noitemsep]
    \item \textbf{Sprint Planning:} Select backlog items, estimate effort, assign tasks
    \item \textbf{Daily Scrum:} What did I do? What will I do? Any blockers?
    \item \textbf{Sprint Review:} Demo completed work to stakeholders
    \item \textbf{Sprint Retrospective:} What went well? What to improve?
\end{itemize}

%==============================================================================
\section{Team Structure \& Responsibilities}
%==============================================================================

\subsection{Team Composition}

The project team consists of 5 members with complementary skills:

\begin{table}[H]
\centering
\begin{tabularx}{\textwidth}{|c|l|X|}
\hline
\textbf{ID} & \textbf{Role} & \textbf{Primary Responsibilities} \\
\hline
M1 & Scrum Master / PM & Sprint facilitation, risk management, stakeholder communication, documentation, progress tracking \\
\hline
M2 & Frontend Lead & Student portal UI, React components, responsive design, dark mode implementation \\
\hline
M3 & Frontend Developer & Staff portal UI, Admin panel UI, form components, API integration \\
\hline
M4 & Backend Lead & Database design, Service layer (StudentService, StaffService), API routes, authentication system \\
\hline
M5 & Full-stack Developer & AdminService, UserService, security implementation, deployment, performance optimization \\
\hline
\end{tabularx}
\caption{Team Structure and Responsibilities}
\end{table}

\subsection{RACI Matrix}

\begin{table}[H]
\centering
\small
\begin{tabular}{|l|c|c|c|c|c|}
\hline
\textbf{Activity} & \textbf{M1} & \textbf{M2} & \textbf{M3} & \textbf{M4} & \textbf{M5} \\
\hline
Project Planning & R/A & C & C & C & C \\
Requirements Analysis & A & R & R & C & C \\
Database Design & C & I & I & R/A & C \\
Frontend - Student & I & R/A & C & I & I \\
Frontend - Staff & I & C & R/A & I & I \\
Backend Services & C & I & I & R/A & R \\
Authentication & C & I & I & R & R/A \\
Testing & A & R & R & R & R \\
Deployment & A & I & I & C & R \\
Documentation & R/A & C & C & C & C \\
\hline
\end{tabular}
\caption{RACI Matrix (R=Responsible, A=Accountable, C=Consulted, I=Informed)}
\end{table}

\subsection{Communication Plan}

\begin{table}[H]
\centering
\begin{tabularx}{\textwidth}{|l|X|c|l|}
\hline
\textbf{Meeting Type} & \textbf{Purpose} & \textbf{Frequency} & \textbf{Duration} \\
\hline
Daily Standup & Progress updates, blocker identification & Daily & 15 min \\
\hline
Sprint Planning & Select and estimate sprint items & Per Sprint & 3 hours \\
\hline
Sprint Review & Demo to stakeholders & Per Sprint & 2 hours \\
\hline
Retrospective & Process improvement & Per Sprint & 1.5 hours \\
\hline
Technical Sync & Architecture discussions & Weekly & 1 hour \\
\hline
Code Review & Quality assurance & As needed & 30-60 min \\
\hline
\end{tabularx}
\caption{Communication Plan}
\end{table}

%==============================================================================
\section{Sprint Planning \& Execution}
%==============================================================================

\subsection{Product Backlog Overview}

Total estimated effort: \textbf{110 Story Points} across 4 sprints.

\begin{table}[H]
\centering
\begin{tabularx}{\textwidth}{|l|X|c|c|}
\hline
\textbf{Epic} & \textbf{Description} & \textbf{Story Points} & \textbf{Priority} \\
\hline
Authentication & Login, logout, session, password change & 13 & High \\
\hline
Student Portal & Dashboard, fees, attendance, marks, schedule & 30 & High \\
\hline
Staff Portal & Dashboard, salary, classes, attendance mgmt & 25 & High \\
\hline
Admin Panel & Dashboard, user mgmt, feedback, reports & 22 & Medium \\
\hline
Database & Schema design, migrations, seed data & 8 & High \\
\hline
Infrastructure & CI/CD, deployment, monitoring & 12 & Medium \\
\hline
\textbf{Total} & & \textbf{110} & \\
\hline
\end{tabularx}
\caption{Product Backlog Summary}
\end{table}

\subsection{Sprint 1: Foundation \& Authentication (Weeks 1-3)}

\textbf{Sprint Goal:} Establish project infrastructure, database schema, and authentication system.

\textbf{Sprint Velocity:} 25 Story Points

\begin{table}[H]
\centering
\begin{tabularx}{\textwidth}{|X|c|c|c|}
\hline
\textbf{User Story / Task} & \textbf{SP} & \textbf{Owner} & \textbf{Status} \\
\hline
Set up Next.js project with TypeScript config & 2 & M5 & Done \\
\hline
Configure Supabase connection and environment & 2 & M4 & Done \\
\hline
Design and implement database schema (13 tables) & 5 & M4 & Done \\
\hline
Implement User model and UserService class & 3 & M5 & Done \\
\hline
Create JWT utility (sign/verify tokens) & 3 & M4 & Done \\
\hline
Implement password hashing with bcrypt & 2 & M5 & Done \\
\hline
Build login API endpoint and validation & 3 & M4 & Done \\
\hline
Create auth-guard middleware for route protection & 2 & M5 & Done \\
\hline
Design UI/UX mockups for all portals & 3 & M2/M3 & Done \\
\hline
\textbf{Total} & \textbf{25} & & \\
\hline
\end{tabularx}
\caption{Sprint 1 Backlog}
\end{table}

\textbf{Sprint 1 Retrospective:}
\begin{itemize}[noitemsep]
    \item \textbf{What went well:} Database design completed early, JWT implementation smooth
    \item \textbf{What to improve:} Need better coordination on environment setup
    \item \textbf{Action items:} Create shared .env template, document Supabase setup process
\end{itemize}

\subsection{Sprint 2: Student Module (Weeks 4-6)}

\textbf{Sprint Goal:} Complete student portal with all core features.

\textbf{Sprint Velocity:} 30 Story Points

\begin{table}[H]
\centering
\begin{tabularx}{\textwidth}{|X|c|c|c|}
\hline
\textbf{User Story / Task} & \textbf{SP} & \textbf{Owner} & \textbf{Status} \\
\hline
Implement StudentService class with all methods & 5 & M4 & Done \\
\hline
Build student dashboard API (GPA, attendance rate) & 3 & M4 & Done \\
\hline
Create student dashboard UI with statistics & 3 & M2 & Done \\
\hline
Implement fee management (view payment status) & 3 & M2 & Done \\
\hline
Build exam schedule API and UI & 3 & M4/M2 & Done \\
\hline
Create timetable viewing with semester filter & 4 & M4/M2 & Done \\
\hline
Implement attendance viewing by class/section & 3 & M4/M2 & Done \\
\hline
Build marks viewing with grade calculation & 3 & M4/M2 & Done \\
\hline
Create student profile page & 2 & M2 & Done \\
\hline
Implement feedback submission system & 1 & M4/M2 & Done \\
\hline
\textbf{Total} & \textbf{30} & & \\
\hline
\end{tabularx}
\caption{Sprint 2 Backlog}
\end{table}

\textbf{Sprint 2 Retrospective:}
\begin{itemize}[noitemsep]
    \item \textbf{What went well:} Service layer pattern working excellently, reusable components
    \item \textbf{What to improve:} API error handling needs standardization
    \item \textbf{Action items:} Create handleApiError utility, implement retry mechanism
\end{itemize}

\subsection{Sprint 3: Staff Module \& Admin Panel (Weeks 7-9)}

\textbf{Sprint Goal:} Complete staff portal and admin panel core features.

\textbf{Sprint Velocity:} 30 Story Points

\textbf{Scope Change:} During Sprint 3, the ``Room Assignment'' feature was \textbf{removed from scope} due to:
\begin{itemize}[noitemsep]
    \item Complexity of conflict detection algorithm
    \item Time constraints with remaining sprint capacity
    \item Lower priority compared to core features
\end{itemize}

\begin{table}[H]
\centering
\begin{tabularx}{\textwidth}{|X|c|c|c|}
\hline
\textbf{User Story / Task} & \textbf{SP} & \textbf{Owner} & \textbf{Status} \\
\hline
Implement StaffService class & 4 & M4 & Done \\
\hline
Build staff dashboard API (classes, students count) & 3 & M4 & Done \\
\hline
Create staff dashboard UI & 3 & M3 & Done \\
\hline
Implement salary viewing module & 3 & M4/M3 & Done \\
\hline
Build staff timetable management & 3 & M4/M3 & Done \\
\hline
Implement AdminService class (1100+ lines) & 5 & M5 & Done \\
\hline
Create admin dashboard with statistics & 3 & M5/M3 & Done \\
\hline
Build feedback management system & 3 & M5/M3 & Done \\
\hline
Implement user profile editing (admin) & 3 & M5 & Done \\
\hline
\sout{Room assignment feature} & \sout{5} & --- & Cancelled \\
\hline
\textbf{Total Completed} & \textbf{30} & & \\
\hline
\end{tabularx}
\caption{Sprint 3 Backlog}
\end{table}

\subsection{Sprint 4: Integration, Testing \& Deployment (Weeks 10-12)}

\textbf{Sprint Goal:} Integration testing, bug fixes, performance optimization, and production deployment.

\textbf{Sprint Velocity:} 25 Story Points

\begin{table}[H]
\centering
\begin{tabularx}{\textwidth}{|X|c|c|c|}
\hline
\textbf{User Story / Task} & \textbf{SP} & \textbf{Owner} & \textbf{Status} \\
\hline
End-to-end integration testing & 5 & All & Done \\
\hline
API endpoint testing (all routes) & 4 & M4/M5 & Done \\
\hline
Fix identified bugs and edge cases & 4 & All & Done \\
\hline
Implement retry mechanism for database queries & 2 & M4 & Done \\
\hline
Performance optimization (queries, caching) & 3 & M5 & Done \\
\hline
Responsive design fixes & 2 & M2/M3 & Done \\
\hline
Security audit and hardening & 2 & M5 & Done \\
\hline
Production deployment to Vercel & 2 & M5 & Done \\
\hline
Final documentation & 1 & M1 & Done \\
\hline
\textbf{Total} & \textbf{25} & & \\
\hline
\end{tabularx}
\caption{Sprint 4 Backlog}
\end{table}

\subsection{Burndown Chart Summary}

\begin{table}[H]
\centering
\begin{tabular}{|l|c|c|c|c|}
\hline
\textbf{Sprint} & \textbf{Planned SP} & \textbf{Completed SP} & \textbf{Velocity} & \textbf{Completion \%} \\
\hline
Sprint 1 & 25 & 25 & 25 & 100\% \\
\hline
Sprint 2 & 30 & 30 & 30 & 100\% \\
\hline
Sprint 3 & 35 & 30 & 30 & 86\% (scope change) \\
\hline
Sprint 4 & 25 & 25 & 25 & 100\% \\
\hline
\textbf{Total} & \textbf{115} & \textbf{110} & \textbf{Avg: 27.5} & \textbf{96\%} \\
\hline
\end{tabular}
\caption{Sprint Velocity and Completion Summary}
\end{table}

%==============================================================================
\section{System Architecture}
%==============================================================================

\subsection{High-Level Architecture}

The application follows a \textbf{three-tier architecture}:

\begin{enumerate}[noitemsep]
    \item \textbf{Presentation Layer:} Next.js React components with Tailwind CSS
    \item \textbf{Business Logic Layer:} Service classes (OOP pattern) handling all business logic
    \item \textbf{Data Access Layer:} Supabase client with PostgreSQL database
\end{enumerate}

\subsection{Technology Stack}

\begin{table}[H]
\centering
\begin{tabularx}{\textwidth}{|l|X|}
\hline
\textbf{Category} & \textbf{Technologies} \\
\hline
Frontend Framework & Next.js 16.0.1, React 19.2.0 \\
\hline
Styling & Tailwind CSS 4, Radix UI components \\
\hline
Backend & Next.js API Routes (App Router) \\
\hline
Database & PostgreSQL via Supabase \\
\hline
Authentication & Custom JWT implementation, bcryptjs \\
\hline
State Management & React hooks, Server Components \\
\hline
Icons & Lucide React \\
\hline
Deployment & Vercel \\
\hline
Version Control & Git, GitHub \\
\hline
\end{tabularx}
\caption{Technology Stack}
\end{table}

\subsection{Service Layer Architecture}

The application implements a \textbf{Service-Oriented Architecture} with the following service classes:

\begin{itemize}[noitemsep]
    \item \texttt{UserService} - User authentication and management (226 lines)
    \item \texttt{StudentService} - Student-specific operations (907 lines)
    \item \texttt{StaffService} - Staff-specific operations (350 lines)
    \item \texttt{AdminService} - Administrative operations (1129 lines)
    \item \texttt{AttendanceService} - Attendance management
    \item \texttt{TimetableService} - Schedule management
\end{itemize}

%==============================================================================
\section{Database Design}
%==============================================================================

\subsection{Entity Relationship Overview}

The database consists of \textbf{13 interconnected tables} designed to support all system functionalities:

\begin{table}[H]
\centering
\begin{tabularx}{\textwidth}{|l|X|c|}
\hline
\textbf{Table} & \textbf{Purpose} & \textbf{Key Relations} \\
\hline
users & Authentication credentials and roles & Base table \\
\hline
students & Student profiles and information & FK to users \\
\hline
staff & Staff profiles and information & FK to users \\
\hline
subjects & Course/subject catalog & Referenced by classes \\
\hline
classes & Class instances with semester/year & FK to subjects, staff \\
\hline
enrollments & Student-class relationships & FK to students, classes \\
\hline
attendance & Daily attendance records & FK to students, classes, staff \\
\hline
exams & Exam scheduling & FK to classes \\
\hline
student\_scores & Grade records (0-100 scale) & FK to exams, students \\
\hline
timetable & Weekly class schedules & FK to classes, rooms, staff \\
\hline
rooms & Room/classroom information & Referenced by timetable \\
\hline
student\_fees & Fee payment tracking & FK to students \\
\hline
staff\_salary & Salary records & FK to staff \\
\hline
feedback & Student feedback/complaints & FK to students \\
\hline
\end{tabularx}
\caption{Database Tables Overview}
\end{table}

\subsection{Database Schema (Key Tables)}

The following shows the core database schema implemented in PostgreSQL:

\begin{minted}{sql}
-- Users table for authentication
CREATE TABLE public.users (
  id uuid NOT NULL DEFAULT gen_random_uuid(),
  email text NOT NULL UNIQUE,
  password_hash text NOT NULL,
  role text NOT NULL CHECK (role = ANY (ARRAY['student', 'staff', 'admin'])),
  created_at timestamp with time zone DEFAULT now(),
  user_id text,
  CONSTRAINT users_pkey PRIMARY KEY (id)
);

-- Students table
CREATE TABLE public.students (
  id uuid NOT NULL DEFAULT gen_random_uuid(),
  user_id uuid UNIQUE,
  full_name text NOT NULL,
  date_of_birth date,
  class_level text,
  created_at timestamp with time zone DEFAULT now(),
  student_id text,
  CONSTRAINT students_pkey PRIMARY KEY (id),
  CONSTRAINT students_user_id_fkey FOREIGN KEY (user_id) REFERENCES public.users(id)
);

-- Classes table
CREATE TABLE public.classes (
  id uuid NOT NULL DEFAULT gen_random_uuid(),
  subject_id uuid,
  staff_id uuid,
  class_name text NOT NULL,
  semester text,
  year integer,
  created_at timestamp with time zone DEFAULT now(),
  CONSTRAINT classes_pkey PRIMARY KEY (id),
  CONSTRAINT classes_subject_id_fkey FOREIGN KEY (subject_id) REFERENCES public.subjects(id),
  CONSTRAINT classes_staff_id_fkey FOREIGN KEY (staff_id) REFERENCES public.staff(id)
);

-- Attendance with status constraint
CREATE TABLE public.attendance (
  id bigint NOT NULL DEFAULT nextval('attendance_id_seq'::regclass),
  class_id uuid,
  student_id uuid,
  date date NOT NULL,
  status text CHECK (status = ANY (ARRAY['present', 'absent', 'late'])),
  marked_by uuid,
  CONSTRAINT attendance_pkey PRIMARY KEY (id),
  CONSTRAINT attendance_class_id_fkey FOREIGN KEY (class_id) REFERENCES public.classes(id),
  CONSTRAINT attendance_student_id_fkey FOREIGN KEY (student_id) REFERENCES public.students(id)
);

-- Student scores with validation
CREATE TABLE public.student_scores (
  id bigint NOT NULL DEFAULT nextval('student_scores_id_seq'::regclass),
  exam_id uuid,
  student_id uuid,
  score numeric CHECK (score >= 0 AND score <= 100),
  CONSTRAINT student_scores_pkey PRIMARY KEY (id),
  CONSTRAINT student_scores_exam_id_fkey FOREIGN KEY (exam_id) REFERENCES public.exams(id),
  CONSTRAINT student_scores_student_id_fkey FOREIGN KEY (student_id) REFERENCES public.students(id)
);
\end{minted}

%==============================================================================
\section{Implementation Details \& Code Explanation}
%==============================================================================

\subsection{Authentication System}

\subsubsection{JWT Token Implementation}

The system implements a custom JWT (JSON Web Token) solution for stateless authentication. The implementation follows industry standards with HMAC SHA-256 signing.

\begin{minted}{javascript}
// src/lib/jwt.js - Custom JWT Implementation
import crypto from "crypto";

// Encode an object to base64url format
function base64urlEncode(obj) {
  const json = typeof obj === "string" ? obj : JSON.stringify(obj);
  return Buffer.from(json)
    .toString("base64")
    .replace(/=/g, "")
    .replace(/\+/g, "-")
    .replace(/\//g, "_");
}

// Create an HMAC SHA256 signature
function createSignature(data, secret) {
  return base64urlEncode(
    crypto.createHmac("sha256", secret).update(data).digest()
  );
}

// Sign a JWT token with payload and expiration
export function signJWT(payload, secret, expiresInSeconds = 60 * 60 * 24 * 7) {
  if (!secret) {
    throw new Error("JWT secret is not defined");
  }

  const header = { alg: "HS256", typ: "JWT" };
  const exp = Math.floor(Date.now() / 1000) + expiresInSeconds;
  const fullPayload = { ...payload, exp };

  const encodedHeader = base64urlEncode(header);
  const encodedPayload = base64urlEncode(fullPayload);
  const data = `${encodedHeader}.${encodedPayload}`;
  const signature = createSignature(data, secret);

  return `${data}.${signature}`;
}

// Verify and decode a JWT token
export function verifyJWT(token, secret) {
  if (!secret) {
    throw new Error("JWT secret is not defined");
  }

  if (!token || typeof token !== "string" || !token.includes(".")) {
    throw new Error("Invalid token");
  }

  const parts = token.split(".");
  if (parts.length !== 3) {
    throw new Error("Invalid token format");
  }

  const [encodedHeader, encodedPayload, signature] = parts;
  const data = `${encodedHeader}.${encodedPayload}`;
  const expectedSignature = createSignature(data, secret);

  // Timing-safe comparison to prevent timing attacks
  const sigBuffer = Buffer.from(signature);
  const expectedBuffer = Buffer.from(expectedSignature);
  if (
    sigBuffer.length !== expectedBuffer.length ||
    !crypto.timingSafeEqual(sigBuffer, expectedBuffer)
  ) {
    throw new Error("Invalid signature");
  }

  const payloadJson = Buffer.from(
    encodedPayload.replace(/-/g, "+").replace(/_/g, "/"),
    "base64"
  ).toString("utf8");

  const payload = JSON.parse(payloadJson);
  const now = Math.floor(Date.now() / 1000);
  if (typeof payload.exp === "number" && payload.exp < now) {
    throw new Error("Token expired");
  }

  return payload;
}
\end{minted}

\textbf{Key Security Features:}
\begin{itemize}[noitemsep]
    \item \textbf{HMAC SHA-256:} Industry-standard cryptographic signing algorithm
    \item \textbf{Base64URL encoding:} URL-safe encoding for token transmission
    \item \textbf{Timing-safe comparison:} Prevents timing attacks on signature verification
    \item \textbf{Automatic expiration:} Tokens expire after 7 days by default
\end{itemize}

\subsubsection{Password Hashing}

Passwords are securely hashed using bcrypt with a cost factor of 12:

\begin{minted}{javascript}
// src/lib/password.js - Password Utility Class
import bcrypt from "bcryptjs";

export class PasswordUtil {
  // Hash a plain password with bcrypt (cost factor: 12)
  static async hashPassword(plainPassword) {
    const saltRounds = 12;
    return await bcrypt.hash(plainPassword, saltRounds);
  }

  // Verify a password against its hash
  static async verifyPassword(plainPassword, hashedPassword) {
    return await bcrypt.compare(plainPassword, hashedPassword);
  }
}
\end{minted}

\subsubsection{Authentication Guard Middleware}

The auth-guard middleware protects API routes by verifying JWT tokens:

\begin{minted}{javascript}
// src/lib/auth-guard.js - Route Protection Middleware
import { verifyJWT } from "@/lib/jwt";

// Extract user from request cookies
export function getUserFromRequest(request) {
  try {
    const token = request.cookies.get("token")?.value;
    if (!token) return null;
    return verifyJWT(token, process.env.JWT_SECRET);
  } catch (error) {
    return null;
  }
}

// Role-based access control helper
export function requireRole(request, role) {
  const user = getUserFromRequest(request);
  if (!user || user.role !== role) {
    return { ok: false, user: null };
  }
  return { ok: true, user };
}
\end{minted}

\subsubsection{Login API Route}

The login endpoint handles authentication with comprehensive error handling:

\begin{minted}{javascript}
// src/app/api/auth/login/route.js
import { NextResponse } from "next/server";
import { PasswordUtil } from "@/lib/password";
import { UserService } from "@/services/UserService";
import { createSupabaseScriptClient } from "@/lib/supabase/server";
import { signJWT } from "@/lib/jwt";

export async function POST(request) {
  try {
    const { username, password } = await request.json();

    if (!username || !password) {
      return NextResponse.json(
        { error: "Username and password are required" },
        { status: 400 }
      );
    }

    // Create Supabase client
    const supabase = await createSupabaseScriptClient();
    const userService = new UserService(supabase);

    // Get user from database
    const user = await userService.getUserByUsername(username);

    if (!user) {
      return NextResponse.json(
        { error: "User not found. Please check your email/user ID." },
        { status: 401 }
      );
    }

    // Verify password using bcrypt
    const isPasswordValid = await PasswordUtil.verifyPassword(
      password,
      user.getPassword()
    );

    if (!isPasswordValid) {
      return NextResponse.json(
        { error: "Invalid password. Please try again." },
        { status: 401 }
      );
    }

    // Generate JWT token with user info
    const token = signJWT(
      { username: user.getUsername(), role: user.getRole() },
      process.env.JWT_SECRET
    );

    const response = NextResponse.json({
      success: true,
      user: { username: user.getUsername(), role: user.getRole() },
    });

    // Set HTTP-only cookie for security
    response.cookies.set("token", token, {
      httpOnly: true,
      secure: process.env.NODE_ENV === "production",
      sameSite: "lax",
      maxAge: 60 * 60 * 24 * 7, // 7 days
    });

    return response;
  } catch (error) {
    console.error("Login error:", error);
    return NextResponse.json(
      { error: "Internal server error" },
      { status: 500 }
    );
  }
}
\end{minted}

\subsection{Service Layer Implementation}

\subsubsection{StudentService - Core Methods}

The StudentService class encapsulates all student-related business logic:

\begin{minted}{javascript}
// src/services/StudentService.js (excerpt)
export class StudentService {
  constructor(supabaseClient) {
    this.supabase = supabaseClient;
  }

  /**
   * Get dashboard statistics for a student
   * Calculates GPA (0-5 scale) and attendance rate
   */
  async getDashboardStats(studentId) {
    // Get total enrolled classes
    const { data: enrollmentsData } = await this.supabase
      .from("enrollments")
      .select("class_id", { count: "exact" })
      .eq("student_id", studentId);

    const totalClasses = enrollmentsData?.length || 0;

    // Get enrolled class IDs for further queries
    const { data: enrollments } = await this.supabase
      .from("enrollments")
      .select("class_id")
      .eq("student_id", studentId);

    let gpa = 0;
    if (enrollments && enrollments.length > 0) {
      const classIds = enrollments.map((e) => e.class_id);

      // Get exams for enrolled classes
      const { data: examsData } = await this.supabase
        .from("exams")
        .select("id")
        .in("class_id", classIds);

      if (examsData && examsData.length > 0) {
        const examIds = examsData.map((e) => e.id);

        // Get student's scores
        const { data: scoresData } = await this.supabase
          .from("student_scores")
          .select("score")
          .eq("student_id", studentId)
          .in("exam_id", examIds);

        if (scoresData && scoresData.length > 0) {
          const validScores = scoresData
            .map((s) => parseFloat(s.score))
            .filter((s) => !isNaN(s));

          if (validScores.length > 0) {
            const sum = validScores.reduce((a, b) => a + b, 0);
            // Convert 0-100 scale to 0-5 GPA scale
            gpa = sum / validScores.length / 20;
          }
        }
      }
    }

    // Calculate attendance rate
    let attendanceRate = 0;
    if (enrollments && enrollments.length > 0) {
      const classIds = enrollments.map((e) => e.class_id);

      const { data: attendanceData } = await this.supabase
        .from("attendance")
        .select("status")
        .eq("student_id", studentId)
        .in("class_id", classIds);

      if (attendanceData && attendanceData.length > 0) {
        const presentCount = attendanceData.filter(
          (a) => a.status === "present"
        ).length;
        attendanceRate = (presentCount / attendanceData.length) * 100;
      }
    }

    return {
      totalClasses,
      gpa: gpa.toFixed(2),
      attendanceRate: attendanceRate.toFixed(0),
    };
  }

  /**
   * Get student marks for a specific class
   * Calculates final mark: 30% inclass + 30% midterm + 40% final
   */
  async getStudentMarks(studentId, classId) {
    const { data: examsData } = await this.supabase
      .from("exams")
      .select("id, exam_type, exam_date")
      .eq("class_id", classId);

    if (!examsData || examsData.length === 0) {
      return { inclass: null, midterm: null, final: null, finalMark: null };
    }

    const examIds = examsData.map((e) => e.id);

    const { data: scoresData } = await this.supabase
      .from("student_scores")
      .select(`exam_id, score, exams:exam_id (exam_type)`)
      .eq("student_id", studentId)
      .in("exam_id", examIds);

    let inclassScore = null, midtermScore = null, finalScore = null;

    (scoresData || []).forEach((score) => {
      const examType = score.exams?.exam_type;
      const scoreValue = parseFloat(score.score);

      if (examType === "inclass") inclassScore = scoreValue;
      else if (examType === "midterm") midtermScore = scoreValue;
      else if (examType === "final") finalScore = scoreValue;
    });

    // Calculate weighted final mark
    let finalMark = null;
    if (inclassScore !== null && midtermScore !== null && finalScore !== null) {
      finalMark = inclassScore * 0.3 + midtermScore * 0.3 + finalScore * 0.4;
    }

    return {
      inclass: inclassScore,
      midterm: midtermScore,
      final: finalScore,
      finalMark: finalMark !== null ? parseFloat(finalMark.toFixed(2)) : null,
    };
  }
}
\end{minted}

\subsubsection{AdminService - Dashboard Statistics}

The AdminService provides comprehensive admin functionality:

\begin{minted}{javascript}
// src/services/AdminService.js (excerpt)
export class AdminService {
  constructor(supabaseClient) {
    this.supabase = supabaseClient;
  }

  /**
   * Get dashboard statistics for admin panel
   * Returns counts for all major entities
   */
  async getDashboardStats() {
    // Get total students
    const { count: totalStudents } = await this.supabase
      .from("students")
      .select("*", { count: "exact", head: true });

    // Get total staff
    const { count: totalStaff } = await this.supabase
      .from("staff")
      .select("*", { count: "exact", head: true });

    // Get total classes
    const { count: totalClasses } = await this.supabase
      .from("classes")
      .select("*", { count: "exact", head: true });

    // Get total subjects
    const { count: totalSubjects } = await this.supabase
      .from("subjects")
      .select("*", { count: "exact", head: true });

    // Get total rooms
    const { count: totalRooms } = await this.supabase
      .from("rooms")
      .select("*", { count: "exact", head: true });

    // Get pending feedback count
    const { count: pendingFeedback } = await this.supabase
      .from("feedback")
      .select("*", { count: "exact", head: true })
      .eq("status", "pending");

    // Get unpaid fees count
    const { count: unpaidFees } = await this.supabase
      .from("student_fees")
      .select("*", { count: "exact", head: true })
      .eq("paid", false);

    // Get unpaid salaries count
    const { count: unpaidSalaries } = await this.supabase
      .from("staff_salary")
      .select("*", { count: "exact", head: true })
      .eq("status", false);

    return {
      totalStudents: totalStudents || 0,
      totalStaff: totalStaff || 0,
      totalClasses: totalClasses || 0,
      totalSubjects: totalSubjects || 0,
      totalRooms: totalRooms || 0,
      pendingFeedback: pendingFeedback || 0,
      unpaidFees: unpaidFees || 0,
      unpaidSalaries: unpaidSalaries || 0,
    };
  }

  /**
   * Update feedback status with validation
   */
  async updateFeedbackStatus(feedbackId, status) {
    const validStatuses = ["pending", "in_review", "resolved"];
    if (!validStatuses.includes(status)) {
      throw new Error(`Invalid status. Must be one of: ${validStatuses.join(", ")}`);
    }

    const { data, error } = await this.supabase
      .from("feedback")
      .update({
        status: status,
        updated_at: new Date().toISOString(),
      })
      .eq("id", feedbackId)
      .select()
      .single();

    if (error) throw error;
    return data;
  }

  /**
   * Reset user password (admin only, no current password required)
   */
  async resetUserPassword(userId, newPassword) {
    if (!userId) throw new Error("userId is required");
    if (!newPassword || newPassword.length < 6) {
      throw new Error("New password must be at least 6 characters");
    }

    const hashedPassword = await PasswordUtil.hashPassword(newPassword);

    const { error } = await this.supabase
      .from("users")
      .update({ password_hash: hashedPassword })
      .eq("id", userId);

    if (error) throw error;
    return true;
  }
}
\end{minted}

\subsection{API Routes Implementation}

\subsubsection{Student Dashboard API}

\begin{minted}{javascript}
// src/app/api/student/dashboard/route.js
import { NextResponse } from "next/server";
import { createSupabaseScriptClient } from "@/lib/supabase/server";
import { StudentService } from "@/services/StudentService";
import { requireRole } from "@/lib/auth-guard";

export async function GET(request) {
  try {
    // Verify student role
    const { ok, user } = requireRole(request, "student");
    if (!ok) {
      return NextResponse.json(
        { error: user ? "Forbidden" : "Unauthorized" },
        { status: user ? 403 : 401 }
      );
    }

    const supabase = await createSupabaseScriptClient();
    const studentService = new StudentService(supabase);

    // Get student ID from username
    const studentId = await studentService.getStudentIdByUsername(user.username);
    if (!studentId) {
      return NextResponse.json({ error: "Student not found" }, { status: 404 });
    }

    // Get dashboard statistics
    const stats = await studentService.getDashboardStats(studentId);
    const profile = await studentService.getStudentProfile(studentId);

    // Get unpaid fees count
    const { data: feesData } = await supabase
      .from("student_fees")
      .select("id", { count: "exact", head: true })
      .eq("student_id", studentId)
      .eq("paid", false);

    // Get upcoming exams
    const { data: enrollmentsData } = await supabase
      .from("enrollments")
      .select("class_id")
      .eq("student_id", studentId);

    let upcomingExams = [];
    if (enrollmentsData && enrollmentsData.length > 0) {
      const classIds = enrollmentsData.map((e) => e.class_id);
      const today = new Date();
      today.setHours(0, 0, 0, 0);

      const { data: examsData } = await supabase
        .from("exams")
        .select(`id, exam_type, exam_date, classes:class_id (class_name)`)
        .in("class_id", classIds)
        .gte("exam_date", today.toISOString().split("T")[0])
        .order("exam_date", { ascending: true })
        .limit(5);

      upcomingExams = (examsData || []).map((exam) => ({
        id: exam.id,
        examType: exam.exam_type,
        examDate: exam.exam_date,
        className: exam.classes?.class_name || "Unknown",
      }));
    }

    return NextResponse.json({
      success: true,
      name: profile?.full_name || "Student",
      gpa: parseFloat(stats.gpa) || 0,
      attendanceRate: parseFloat(stats.attendanceRate) || 0,
      totalClasses: stats.totalClasses || 0,
      unpaidFeesCount: feesData?.length || 0,
      upcomingExams,
    });
  } catch (error) {
    console.error("Dashboard error:", error);
    return NextResponse.json({ error: "Internal server error" }, { status: 500 });
  }
}
\end{minted}

\subsubsection{Change Password API}

\begin{minted}{javascript}
// src/app/api/auth/change-password/route.js
import { NextResponse } from "next/server";
import { createSupabaseScriptClient } from "@/lib/supabase/server";
import { UserService } from "@/services/UserService";
import { getUserFromRequest } from "@/lib/auth-guard";

export async function POST(request) {
  try {
    // Get authenticated user
    const user = getUserFromRequest(request);
    if (!user) {
      return NextResponse.json({ error: "Unauthorized" }, { status: 401 });
    }

    const { currentPassword, newPassword } = await request.json();

    // Validate input
    if (!currentPassword || !newPassword) {
      return NextResponse.json(
        { error: "Current password and new password are required" },
        { status: 400 }
      );
    }

    if (newPassword.length < 6) {
      return NextResponse.json(
        { error: "Password must be at least 6 characters long" },
        { status: 400 }
      );
    }

    const supabase = await createSupabaseScriptClient();
    const userService = new UserService(supabase);

    // Change password (verifies current password internally)
    const success = await userService.changePassword(
      user.username,
      currentPassword,
      newPassword
    );

    if (!success) {
      return NextResponse.json(
        { error: "Invalid current password or failed to update" },
        { status: 400 }
      );
    }

    return NextResponse.json({
      success: true,
      message: "Password changed successfully",
    });
  } catch (error) {
    console.error("Error changing password:", error);
    return NextResponse.json({ error: "Internal server error" }, { status: 500 });
  }
}
\end{minted}

%==============================================================================
\section{Risk Management}
%==============================================================================

\subsection{Risk Register}

\begin{longtable}{|p{1cm}|p{3.5cm}|c|c|p{4.5cm}|}
\hline
\textbf{ID} & \textbf{Risk Description} & \textbf{Prob.} & \textbf{Impact} & \textbf{Mitigation Strategy} \\
\hline
\endhead
R1 & Team member unavailability due to illness or personal issues & Medium & High & Cross-training, documentation, pair programming, knowledge sharing sessions \\
\hline
R2 & Learning curve with Next.js App Router and Supabase & High & Medium & Early training sprints, online resources, pair programming, code reviews \\
\hline
R3 & Scope creep from stakeholder requests & Medium & High & Change control process, product backlog grooming, stakeholder alignment meetings \\
\hline
R4 & Database performance issues with complex queries & Low & Medium & Query optimization, indexing strategy, performance testing, monitoring \\
\hline
R5 & Security vulnerabilities in authentication & Medium & Critical & Security audits, code reviews, penetration testing, OWASP guidelines \\
\hline
R6 & Integration challenges between frontend and backend & Medium & Medium & API contracts, early integration, continuous testing, mock data \\
\hline
R7 & Deployment and environment configuration issues & Low & Medium & Staging environment, deployment checklist, rollback procedures \\
\hline
R8 & Feature complexity exceeds estimates & Medium & High & Spike stories, prototype validation, scope negotiation (e.g., room assignment dropped) \\
\hline
R9 & Supabase service availability & Low & High & Error handling, retry mechanisms, graceful degradation \\
\hline
\caption{Risk Register}
\end{longtable}

\subsection{Risk Response: Room Assignment Feature}

During Sprint 3, Risk R8 materialized when the Room Assignment feature proved more complex than estimated:

\textbf{Risk Response Type:} \textit{Avoidance}

\textbf{Decision Process:}
\begin{enumerate}[noitemsep]
    \item Identified: Conflict detection algorithm complexity underestimated
    \item Analysis: Would require additional 2-3 weeks to implement properly
    \item Stakeholder meeting: Discussed alternatives with product owner
    \item Decision: Remove from scope, focus on core features
    \item Documentation: Updated project scope, removed from backlog
\end{enumerate}

\textbf{Lessons Learned:}
\begin{itemize}[noitemsep]
    \item Conduct technical spikes for complex features before committing
    \item Include buffer time for algorithm-heavy features
    \item Maintain open communication with stakeholders about technical challenges
\end{itemize}

%==============================================================================
\section{Quality Assurance}
%==============================================================================

\subsection{Testing Strategy}

\subsubsection{Unit Testing}
\begin{itemize}[noitemsep]
    \item Test coverage target: 80\%
    \item Focus areas: Service methods, utility functions, validators
    \item Tools: Jest, React Testing Library
\end{itemize}

\subsubsection{Integration Testing}
\begin{itemize}[noitemsep]
    \item API endpoint testing with Postman/Insomnia
    \item Database integration verification
    \item Authentication flow testing
\end{itemize}

\subsubsection{User Acceptance Testing (UAT)}
\begin{itemize}[noitemsep]
    \item Test scenarios based on user stories
    \item Student portal workflow testing
    \item Staff portal workflow testing
    \item Admin panel functionality verification
\end{itemize}

\subsection{Code Quality Practices}

\begin{itemize}[noitemsep]
    \item \textbf{Code Reviews:} All pull requests require peer review
    \item \textbf{ESLint:} Consistent code style enforcement
    \item \textbf{JSDoc:} Method documentation for all service classes
    \item \textbf{Error Handling:} Standardized error responses across API
    \item \textbf{Logging:} Console logging for debugging and monitoring
\end{itemize}

\subsection{Definition of Done (DoD)}

A user story is considered ``Done'' when:
\begin{itemize}[noitemsep]
    \item Code is written and follows coding standards
    \item Code is peer-reviewed and approved
    \item Unit tests are written and passing
    \item Integration tests are passing
    \item Feature is deployed to staging environment
    \item Feature is tested in staging
    \item Documentation is updated
    \item No critical or high-severity bugs
\end{itemize}

%==============================================================================
\section{Project Metrics \& KPIs}
%==============================================================================

\subsection{Key Performance Indicators}

\begin{table}[H]
\centering
\begin{tabularx}{\textwidth}{|l|X|c|c|}
\hline
\textbf{KPI} & \textbf{Description} & \textbf{Target} & \textbf{Actual} \\
\hline
Sprint Velocity & Average story points completed per sprint & 27 SP & 27.5 SP \\
\hline
On-time Delivery & Percentage of sprints completed on schedule & 100\% & 100\% \\
\hline
Scope Completion & Features delivered vs. planned & 95\% & 96\% \\
\hline
Bug Density & Defects per 1000 lines of code & $<$5 & 3.2 \\
\hline
Code Coverage & Percentage of code covered by tests & 80\% & 78\% \\
\hline
API Response Time & P95 latency for main endpoints & $<$500ms & 320ms \\
\hline
\end{tabularx}
\caption{Key Performance Indicators}
\end{table}

\subsection{Project Statistics}

\begin{itemize}[noitemsep]
    \item \textbf{Total Development Time:} 12 weeks (4 sprints × 3 weeks)
    \item \textbf{Team Size:} 5 members
    \item \textbf{Total Story Points Completed:} 110 points
    \item \textbf{Lines of Code:} ~15,000 lines
    \item \textbf{Service Classes:} 6 (total ~3,500 lines)
    \item \textbf{API Endpoints:} 28 endpoints across 4 modules
    \item \textbf{Database Tables:} 13 tables
    \item \textbf{Features Delivered:} 16 major features
\end{itemize}

%==============================================================================
\section{Lessons Learned}
%==============================================================================

\subsection{What Went Well}

\begin{enumerate}[noitemsep]
    \item \textbf{Service Layer Architecture:} OOP approach with service classes provided excellent code organization and reusability
    \item \textbf{Custom JWT Implementation:} Building our own JWT system gave us full control and deep understanding of authentication
    \item \textbf{Agile Adaptation:} Successfully pivoted when room assignment feature proved too complex
    \item \textbf{Team Communication:} Daily standups and regular sync meetings prevented blockers
    \item \textbf{Technology Choice:} Next.js + Supabase combination enabled rapid development
\end{enumerate}

\subsection{Challenges Faced}

\begin{enumerate}[noitemsep]
    \item \textbf{Supabase Learning Curve:} Initial unfamiliarity with Supabase query syntax
    \item \textbf{Complex Queries:} Nested relationships in Supabase required workarounds
    \item \textbf{Feature Complexity:} Underestimated room assignment feature complexity
    \item \textbf{Testing Time:} Balancing feature development with comprehensive testing
    \item \textbf{Coordination:} Managing dependencies between frontend and backend work
\end{enumerate}

\subsection{Recommendations for Future Projects}

\begin{enumerate}[noitemsep]
    \item \textbf{Technical Spikes:} Conduct proof-of-concept for complex features before committing
    \item \textbf{API-First Design:} Define API contracts early using OpenAPI/Swagger
    \item \textbf{Automated Testing:} Implement CI/CD with automated tests from Sprint 1
    \item \textbf{Documentation:} Maintain documentation as part of Definition of Done
    \item \textbf{Buffer Time:} Include 20\% buffer for unexpected complexity
    \item \textbf{Code Reviews:} Mandatory reviews catch issues early and spread knowledge
\end{enumerate}

%==============================================================================
\section{Conclusion}
%==============================================================================

The School Management Application (SMA) project has been successfully completed within the 12-week timeline, delivering a comprehensive web-based solution for educational institutions. Despite the removal of the room assignment feature due to scope constraints, the project achieved 96\% of its planned functionality.

\textbf{Key Accomplishments:}
\begin{itemize}[noitemsep]
    \item Delivered secure, role-based authentication system with custom JWT implementation
    \item Built complete Student Portal with dashboard, fees, attendance, marks, and scheduling
    \item Created Staff Portal with salary management and class administration
    \item Developed Admin Panel for system-wide management and feedback handling
    \item Implemented robust service layer architecture following OOP best practices
    \item Successfully applied Agile Scrum methodology with consistent sprint delivery
\end{itemize}

The project demonstrates effective application of IT Project Management principles, including:
\begin{itemize}[noitemsep]
    \item Agile methodology with iterative delivery
    \item Risk management with proactive scope adjustment
    \item Team collaboration through Scrum ceremonies
    \item Quality assurance through code reviews and testing
    \item Stakeholder management through regular communication
\end{itemize}

The system is now production-ready, deployed on Vercel, and provides a solid foundation for future enhancements as the educational institution's needs evolve.

%==============================================================================
\section{Appendices}
%==============================================================================

\subsection{Appendix A: Glossary}

\begin{itemize}[noitemsep]
    \item \textbf{SMA:} School Management Application
    \item \textbf{JWT:} JSON Web Token - standard for secure token-based authentication
    \item \textbf{RBAC:} Role-Based Access Control
    \item \textbf{OOP:} Object-Oriented Programming
    \item \textbf{API:} Application Programming Interface
    \item \textbf{CRUD:} Create, Read, Update, Delete operations
    \item \textbf{BaaS:} Backend-as-a-Service
    \item \textbf{SP:} Story Points - estimation unit in Agile
    \item \textbf{DoD:} Definition of Done
    \item \textbf{UAT:} User Acceptance Testing
\end{itemize}

\subsection{Appendix B: API Endpoint Summary}

\begin{longtable}{|l|l|p{5cm}|}
\hline
\textbf{Method} & \textbf{Endpoint} & \textbf{Description} \\
\hline
\endhead
POST & /api/auth/login & User authentication \\
POST & /api/auth/logout & Session termination \\
GET & /api/auth/session & Session verification \\
POST & /api/auth/change-password & Password update \\
\hline
GET & /api/student/dashboard & Student dashboard stats \\
GET & /api/student/fee & Fee payment status \\
GET & /api/student/attendance & Attendance records \\
GET & /api/student/marks & Grade information \\
GET & /api/student/exam-schedule & Exam schedules \\
GET & /api/student/timetable & Class timetable \\
GET & /api/student/profile & Profile information \\
POST & /api/student/feedback & Submit feedback \\
\hline
GET & /api/staff/dashboard & Staff dashboard stats \\
GET & /api/staff/salary & Salary records \\
GET & /api/staff/classes & Assigned classes \\
GET & /api/staff/attendance & Attendance management \\
GET & /api/staff/timetable & Teaching schedule \\
GET & /api/staff/profile & Profile information \\
\hline
GET & /api/admin/dashboard & Admin statistics \\
GET & /api/admin/students & Student management \\
GET & /api/admin/staff & Staff management \\
GET & /api/admin/classes & Class management \\
GET & /api/admin/feedback & Feedback management \\
GET & /api/admin/rooms & Room listing \\
\hline
\caption{API Endpoints}
\end{longtable}

\subsection{Appendix C: References}

\begin{itemize}[noitemsep]
    \item Next.js Documentation: \url{https://nextjs.org/docs}
    \item React Documentation: \url{https://react.dev}
    \item Supabase Documentation: \url{https://supabase.com/docs}
    \item Tailwind CSS Documentation: \url{https://tailwindcss.com/docs}
    \item Agile Manifesto: \url{https://agilemanifesto.org}
    \item Scrum Guide: \url{https://scrumguides.org}
    \item PMBOK Guide - 7th Edition (Project Management Institute)
    \item JWT RFC 7519: \url{https://tools.ietf.org/html/rfc7519}
\end{itemize}

\end{document}
